\begin{abstract}
When IBLT is used for set reconciliation, the capacity must be chosen while the number of differing elements $d$ is
unknown. A self-sizing protocol proceeds as follows: a first round attempts peeling decoding with a small sketch, and
after failure estimates $d$ from the energy of the count array, then sets the second-round capacity to a lower
quantile of the estimate. This quantile must be computed conditionally on failure, because the measurement and the
success of peeling are determined by the same random mapping, so the distribution of the measurement on failing
samples can deviate from the unconditional distribution. We study this deviation in the ideal peeling model: the
number of hash functions $k\ge3$ is fixed, each element independently and uniformly selects $k$ cells, the elements
carry $\pm1$ signs, and the sign composition is fixed and independent of the mapping.

Within the critical window where the failure probability is nondegenerate, we prove that the standardized energy and a
statistic that depends only on the unsigned mapping and governs the success of peeling (the precritical score) have a
joint Gaussian limit, with correlation coefficient
$\eta_k\rho^2$, where $\rho$ is the sign imbalance and $\eta_k$ has a closed form. Consequently, the
measurement conditional on failure follows the distribution of a bivariate normal truncated on a half-plane, which we
call the Gaussian-selection
law. The coefficient $\eta_k$ changes sign between $k=4$ and $k=5$, and the direction of the shift in the
conditional mean and the lower
quantile after failure reverses accordingly: when $k\ge5$ and the imbalance leans toward one side, keeping the
unconditional threshold
makes the risk of insufficient second-round capacity higher than the nominal level.

The key tool for separating the role of the signs from that of the mapping is a quenched CLT for a signed
quadratic form on a random overlap graph, with the sign composition fixed. As a corollary, when the signs are
asymptotically balanced, conditioning on any unsigned mapping event of probability bounded away from zero and
invariant under element-label permutation does not change the standard normal limit.

After failure the protocol can observe the total sign. Inverting the selection law with it yields an upper confidence
bound for the unknown $d$; this bound is asymptotically conservative within the critical window. As long as
$d/\sqrt M\to\infty$, it extends to the entire load axis. In the numerical experiments, all theoretical predictions
are given by analytic constants, with no parameters fitted. The experiments show that the unconditional threshold
carries a selection error that does not vanish with scale, while the selection-law threshold with a finite-size
correction matches the accuracy of the per-configuration empirical quantile.
\end{abstract}
